%==============================================================================
% Section 3: Mode Dictionary: S^3 x S^1 under EC-Weyl
%==============================================================================
\section{Mode Dictionary: $\Sthree\times\Sone$ under EC--Weyl}
\label{sec:s3}

This section determines the effect of the EC--Weyl coupling on the homogeneous
geometric degrees of freedom of $\Sthree\times\Sone$.  Since the non-propagation
of torsion in the spin-0 sector is already established by Theorem~\ref{thm:palatini},
we focus here on the specific consequences for $\Sthree$: Theorem~4 (Palatini
universality) and its reflection in the mode dictionary.

\subsection{Theorem~4: Palatini Universality}
\label{sec:s3:thm4}

\begin{theorem}[Palatini universality]\label{thm:paluniv}
  In $\Sthree\times\Sone$, the spin-2 and spin-1 masses in the AX and VT
  backgrounds are identical (SymPy exact)\textup{:}

  \textbf{Spin-2 mass (five-fold degenerate)\textup{:}}
  \begin{equation}
    \mu^2_{\mathrm{spin2}}(\Sthree)
    = \frac{128\pi^2 Lr}{3\kappa^2} - \frac{16384\pi^2 L\alpha}{3r}.
  \end{equation}

  \textbf{Spin-1 mass (three-fold degenerate)\textup{:}}
  \begin{equation}
    m^2_{\omega_k}(\Sthree) = \frac{16\pi^2 L^3}{\kappa^2 r},
    \qquad k=0,1,2.
  \end{equation}

  Both expressions are independent of $\eta$ and $V$: the mode spectrum does
  not depend on the type of torsion background.
\end{theorem}

The VT$-$AX difference for spin-2 mass is exactly zero (SymPy exact).  The LC
contribution ($\alpha=0$) to the spin-2 mass is $128\pi^2 Lr/(3\kappa^2)$ and
the EC correction is $-16384\pi^2 L\alpha/(3r)$, the latter depending only on
the topology ($\Sthree$ geometry) and not on $\eta$ or $V$.

\begin{remark}
  The numerical evaluation at $r=3$, $L=\kappa=1$, $\alpha=0$ gives
  $\mu^2_{\mathrm{spin2}} = 128\pi^2 \approx 1263.3$ and
  $m^2_{\omega_k} = 16\pi^2/3 \approx 52.6$.  The LC paper~III uses $r=1$-based
  notation, giving $128\pi^2/3\approx 421.2$ for the spin-2 mass.  The
  discrepancy with the value $48\pi^2$ quoted there arises from the choice of
  radius combined with conventions for the spin-2 basis/Hessian.
\end{remark}

\paragraph{Physical interpretation.}
Palatini universality means that the difference between the AX background
($\eta\neq 0$) and the VT background ($V\neq 0$) has no effect on the
spin-2/spin-1 mass spectrum.  The kinematic part coming from the $\Sthree$
geometry ($C^i_{jk}=4\varepsilon_{ijk}/r$) determines the masses, while
torsion deforms only the shape of the potential (a consequence of Palatini
protection, Theorem~\ref{thm:palatini}).
Script: \texttt{paper03ec/s3\_vt\_spin\_masses.py}.

\subsection{Mode Dictionary Table ($\Sthree\times\Sone$)}
\label{sec:s3:table}

Table~\ref{tab:s3} summarises the mode dictionary of $\Sthree\times\Sone$
under EC--Weyl.

\begin{table}[H]
\centering
\resizebox{\linewidth}{!}{%
\begin{tabular}{lccc}
\toprule
Quantity & AX ($V=0,\,\eta\neq 0$) & VT ($\eta=0,\,V\neq 0$) & MX ($V\neq 0,\,\eta\neq 0$) \\
\midrule
$C^2_\EC$ & $=C^2_\LC=0$ & $=C^2_\LC=0$ & $16V^2\eta^2/(3r^2)$ \\
AX/VT dropout & \checkmark & \checkmark & $\times$ \\
spin-0: $r$ & propagating & propagating & propagating \\
spin-0: $\eta,V$ & $\eta$ non-prop., $V=0$ (Thm.~2) & $\eta=0$, $V$ non-prop.\ (Thm.~2) & both non-prop.\ (Thm.~2) \\
spin-2 mass $\mu^2$ & $128\pi^2 Lr/(3\kappa^2)-16384\pi^2 L\alpha/(3r)$ & same as AX (Thm.~4) & --- \\
spin-1 mass $m^2_\omega$ & $16\pi^2 L^3/(\kappa^2 r)$ (3-fold) & same as AX (Thm.~4) & --- \\
\bottomrule
\end{tabular}%
}
\caption{$\Sthree\times\Sone$ mode dictionary under EC--Weyl.}
\label{tab:s3}
\end{table}

Regarding the MX sector of $\Sthree\times\Sone$: the supplementary script \\
\texttt{paper03ec/s3\_mx\_vacuum\_search.py} checked the stationarity conditions
and Hessian within the range $|\eta|\le 3$.  No stable full stationary vacuum
satisfying both $\nabla_{(r,\eta,V)}\Veff=0$ and positive-definite Hessian was
found.  Real stationary points may appear on the $\alpha>0$ side, but all of
them have at least one negative eigenvalue and are saddle points.  The finite-radius
well appearing in a fixed-$V$ slice is a one-dimensional minimum in the $r$
direction and does not constitute a full stationary vacuum.

\subsection{Pontryagin Structure: AX = Pl\"{u}cker, VT = Pfaffian, MX $\neq 0$}
\label{sec:s3:pontryagin}

The Pontryagin protection structure is inherited from paper~III.

\textbf{AX mode (Pl\"ucker protection):} When $V=0$, the Pontryagin density
$P=\innerprod{R}{\star R}$ vanishes identically by a Pl\"ucker relation.
AX torsion contributes no antisymmetric part to the Riemann tensor that could
activate $P$.

\textbf{VT mode (Pfaffian protection):} When $\eta=0$, $P$ vanishes via a
Pfaffian-type identity.  VT torsion modifies only the symmetric part of the
Riemann tensor, leaving the antisymmetric combination that enters $P$ zero.

\textbf{MX mode:} When $V\neq 0,\eta\neq 0$, the Pontryagin density at the
isotropic point is (from paper~III):
\begin{equation}
  P_0 = \frac{2V\eta(V^2r^2 + 9\eta^2 - 36)}{9r^3}.
\end{equation}

\paragraph{Effect of EC.}
The Weyl coupling $\alpha$ enters only the action, not the definition of the
Pontryagin density $P$ itself.  Therefore the AX/VT/MX Pontryagin classification
is unchanged from paper~III, and what changes in the EC extension is the
effective potential $\Veff$ side.

\subsection{EC-Specific Modification under EC--Weyl}
\label{sec:s3:ec}

\textbf{Preservation of AX/VT stability.}
Under EC correction $\alpha<0$, the structure of the AX/VT stable vacua is
unchanged from LC.  By the AX/VT dropout (Theorem~\ref{thm:dropout}),
$C^2_\EC = C^2_\LC$, so the Weyl coupling does not modify the effective
potential.

\textbf{EC-specific cubic coupling ($\alpha$-cubic).}
As a new nonlinear coupling absent in LC, the following appears:
\begin{equation}
  \frac{\partial^3 V_\EC}{\partial\alpha\,\partial\eta\,\partial V}
  = -\frac{128\pi^2 V\eta r}{3} \neq 0
  \qquad (\Sthree).
\end{equation}
This gives the EC-specific cubic coefficient in the homogeneous EFT with respect
to $\alpha$, a cubic channel absent in LC (script:
\texttt{paper03ec/ec\_cubic\_vertex.py}).  This term does not appear in the
second-order mode dictionary of Table~\ref{tab:s3} and enters only at cubic
order in the homogeneous EFT.  The coefficient corresponds to the
$\Sthree\times\Sone$ volume convention $\Vol(\Sthree\times\Sone)=2\pi^2 Lr^3$
of Appendix~\ref{app:A}.  For comparison, the NY-cubic channel
$\partial^3 V/\partial\theta\partial\eta\partial V = 4\pi^2 r^2$ is invariant
under EC (a consequence of Palatini protection).

In the auxiliary formal MIXING sector, on the other hand,
\begin{equation}
  C_\delta(\Sthree) = 0,
  \qquad
  \frac{\partial C_\delta(\Sthree)}{\partial\alpha} = 0
\end{equation}
(Theorem~\ref{thm:delta}, Appendix~\ref{app:D}).  Therefore the new homogeneous
cubic structure activated by EC extension is confined to the physical $\eta$-$V$
channel, and the $\delta$-sector remains null.
